\centerline{\bf \Huge Оценка + Пример II}

\begin{flushright}
        \scriptsize{Твои возможности обычно больше, чем твоя их оценка\\ Шри Ауробиндо}
\end{flushright}


\problem{\textbf{1}} В городе Честервилле солнце светит нечасто: среди любых пяти дней подряд есть хотя бы четыре пасмурных. Зато среди любых шести дней подряд найдётся хотя бы один солнечный. Сколько солнечных дней может быть в Честервилле в сентябре? Укажите все возможные варианты.

\problem{\textbf{2}} Шахматную доску $8\times8$ перекрасили в несколько цветов (каждую клетку – в один цвет). Оказалось, что если две клетки – соседние по диагонали или отстоят друг от друга на ход коня, то они обязательно разного цвета. Какое наименьшее число цветов могло быть использовано?

\problem{\textbf{3}} В некоторых клетках таблицы $10\times10$ расставлены несколько крестиков и несколько ноликов. Известно, что нет линии (строки или столбца), полностью заполненной одинаковыми значками (крестиками или ноликами). Однако, если в любую пустую клетку поставить любой значок, то это условие нарушится. Какое минимальное число значков может стоять в таблице?

\problem{\textbf{4}} Для круглосуточной охраны объекта нужно установить дежурство на посту в две смены: дневную и ночную. Дежурный может отработать дневную или ночную смену, или же сутки подряд. В первом случае сразу после дежурства ему предоставляется отдых не менее одних суток, во втором — не менее полутора суток, в третьем — не менее 2,5 суток. Какое наименьшее количество дежурных необходимо при этих условиях?


\problem{\textbf{5}} В каждой клетке полоски длины 100 стоит по фишке. Можно за 1 рубль поменять местами любые две соседние фишки, а также можно бесплатно поменять местами любые две фишки, между которыми стоят ровно 4 фишки. За какое наименьшее количество рублей можно переставить фишки в обратном порядке?\\


\centerline{\textit{1 балл за Пример + 1 балл за Оценку}}\\